\section{Physical-Robot Setup and Evaluation}
\label{app:real-robot}

This section documents the physical-robot experiments in Section~\ref{sec:behavior-cloning-robot-control}.

\subsection{Platform and Data Collection}
\label{app:robot-data}

Demonstrations are collected by joint-space teleoperation from a
\methodname{reBot Arm 102} leader to a \methodname{reBot B601-DM} follower through
\methodname{LeRobot}~\citep{cadene2026lerobot}. The follower uses impedance control
with gravity feedforward. One proficient operator collects successful demonstrations
without intentional exploration. Between episodes, both arms return to the calibration
zero pose before recording resumes.

A \methodname{ZED-Mini} records the scene at $672\times376$ pixels per eye, and a
wrist-mounted \methodname{Logitech C270} records at $640\times360$ pixels. The model
uses the left scene view and wrist view, both cropped to 256$\times$256 pixels; the right stereo view is recorded but unused.
At 30 Hz, a Raspberry Pi 5 stores the three video streams, seven-dimensional joint
targets as actions, and follower position, torque, and velocity as a 21-dimensional state.
% TODO: replace the dataset-link placeholder with the public release URL after submission.

\begin{table}[H]
\centering
\small
\caption{Datasets for the physical-robot tasks in Section~\ref{sec:behavior-cloning-robot-control}.}
\label{tab:robot-current-datasets}
\label{tab:robot-datasets}
\begin{tabular}{@{}lrrr>{\raggedright\arraybackslash}p{0.38\linewidth}@{}}
\toprule
Task & Episodes & Frames & Duration & Task Description \\
\midrule
Pick and Place T & 200 & 52,297 & 29 min & "Pick up the T block and place it onto the outline." \\
Push-T & 100 & 61,996 & 34 min & "Push the T block onto the outline." \\
Tape Stack & 100 & 23,726 & 13 min & "Stack the green roll of tape on the white roll of tape." \\
\bottomrule
\end{tabular}
\end{table}

\subsection{Model Implementation}
\label{app:robot-implementation}

The \method{} models in Section~\ref{sec:behavior-cloning-robot-control} use the state MSE and action flow
matching objectives from Section~\ref{sec:joint-predictor}. In the frozen
\methodname{V-JEPA 2.1} variant, attentive pooling and a BatchNorm MLP reduce each
camera frame to one summary token. Proprioceptive state and previous actions are
linearly projected from raw vectors to the model dimension. Learned queries initialize
the future-camera and action streams.

The real world \method{} variants take two video context frames and predict three future frames,
all sampled at 2.5~Hz.
They use model width 512 and depth eight, and train with AdamW
at batch size 256, using gradient accumulation when needed. Both variants use a
ViT-B visual backbone. The end-to-end model trains its visual encoder from scratch;
the frozen variant initializes its encoder from \methodname{V-JEPA 2.1}~\citep{murlabadia2026vjepa21}
and keeps it fixed. The action chunk length and execution interval are specified
in Appendix~\ref{app:bc-physical}.

For Table~\ref{tab:real-robot-results}, batch-1 GFLOPs include all inference steps
for one full action chunk: 100 DDPM steps for \methodname{Diffusion Policy}~\citep{chi2023diffusionpolicy},
ten Euler steps for $\pi_{0.5}$~\citep{physicalintelligence2025pi05}, and ten
Euler steps for each physical-robot \method{} variant.

\subsection{Physical-Robot Baseline Configurations}
\label{app:bc-physical}
For the comparison in Section~\ref{sec:behavior-cloning-robot-control}, both baselines use the
\methodname{LeRobot} implementations~\citep{cadene2026lerobot}, adapted to the
scene and wrist cameras and seven-dimensional joint-target actions. Each policy
returns 36 future actions and executes the first 18 at 30~Hz before requesting
a new chunk, matching the \method{} variants. The returned and executed chunks
span 1.2 and 0.6 seconds, respectively. The demonstrations and evaluation procedure
are described in Appendices~\ref{app:robot-data} and~\ref{app:robot-evaluation}, with 30 trials per task and method.

The \dataset{Pick and Place T} and \dataset{Push-T} training configurations resize both camera
streams to $256\times256$ and retain seven joint torques from the recorded
proprioceptive state. Actions and torques are normalized by mapping their
training-set first and 99th percentiles to $-1$ and $1$, respectively.
Training examples start every three recorded frames; the observation offsets
and action targets retain their 30~Hz timing. We hold out 10\% of episodes
using split seed zero. The configurations specify 20,000 optimizer updates
with effective batch size 256, a 500-update linear learning-rate warmup,
and cosine decay. Validation runs every 200 updates; the trainer saves both
the final checkpoint and the checkpoint with the lowest validation loss.

\baselineitem{MethodDiffusionPolicy}{\textbf{\methodname{Diffusion Policy}}~\citep{chi2023diffusionpolicy}.
The robot implementation uses a separate ImageNet-initialized ResNet-18 with
spatial softmax for each camera and a conditional one-dimensional U-Net with
channel widths $(336,672,1344)$. It conditions on two consecutive image and
torque observations. Images are randomly cropped to $243\times243$ during
training and center-cropped at evaluation. The internal diffusion horizon is
40 action positions, aligned from one step before the current observation;
the policy returns positions corresponding to the next 36 actions. This
40-step internal sequence satisfies the U-Net's temporal downsampling constraint.
Training uses noise-prediction MSE with a squared-cosine schedule, and inference
uses 100 DDPM steps. AdamW uses learning rate $10^{-4}$, betas $(0.95,0.999)$,
and weight decay $10^{-6}$. Exponential moving averaging with decay 0.999
supplies the weights used for validation and saved checkpoints.}

\baselineitem{PaperGray}{\textbf{$\pi_{0.5}$}~\citep{physicalintelligence2025pi05}.
The policy conditions on the current scene and wrist images together with
the task instruction and current torque vector. Its image frontend resizes
the inputs to $224\times224$. The wrapper discretizes the normalized torque
values into 256 bins and includes them in the language prompt.
The action expert predicts a 36-action chunk by flow matching, using ten
Euler steps at inference. The expert-only configuration freezes the vision
and language backbone while training the action expert and action/time
projections. Its task configurations use AdamW with learning rate
$2.5\times10^{-5}$ and weight decay $10^{-4}$, and gradient accumulation to
reach effective batch size 256.}



\subsection{Closed-Loop Evaluation}
\label{app:robot-evaluation}

 For Section~\ref{sec:behavior-cloning-robot-control}, the robot sends its observation window to a
remote inference machine, receives an action chunk, and executes it at 30 Hz under
the receding-horizon schedule. A human evaluator labels each episode as success or
failure, using the same criteria used during teleoperation. Each method is evaluated in 30 trials per task.
% TODO: record per-trial outcomes, confidence intervals, execution
% horizons, and measured model-only and network-inclusive latency for each configuration.

\subsection{Representation Diagnostics and Evaluation}
\label{app:robot-diagnostics}

For visualization, we fit a pixel decoder $d_\psi$ to frozen encoder outputs
$\vz_t=\enc(\vo_t)$. Learnable per-patch queries use cross-attention to map these
representations to pixel patches. Only the decoder parameters $\psi$ are trained,
using per-pixel MSE over training frames:
\begin{equation}
    \mathcal{L}_{\mathrm{dec}}(\psi)
    =\mathbb{E}_{\vo_t}\!\left[
    \operatorname{MSE}\bigl(d_\psi(\enc(\vo_t)),\vo_t\bigr)\right].
    \label{eq:robot-pixel-decoder}
\end{equation}
The encoder and latent dynamics remain frozen. Framewise reconstructions are
$d_\psi(\vz_t)$, and decoded rollout predictions are $d_\psi(\hat{\vz}_t)$.

Figure~\ref{fig:robot-rollout-scene} compares both encoder variants on the same
held-out Pick and Place T episode from the scene camera. Each strip shows recorded
video, framewise reconstructions, and decoded autoregressive predictions. Latent
rollouts start from the initial observation context and advance under recorded
actions, without intermediate ground-truth frames or newly generated actions.
We display ten frames at indices 0, 3, 6, 9, 12, 15, 18, 21, 24, and 26, spanning
10.4 seconds at the 2.5~Hz model sampling rate.

\begin{figure}[htbp!]
    \centering
    {\small\bfseries End-to-end LeWAM\par}
    \includegraphics[width=\textwidth]{figures/real_robot_rollouts/output/lewam_scene.pdf}
    \vspace{2pt}
    {\small\bfseries LeWAM with frozen \methodname{V-JEPA 2.1}\par}
    \includegraphics[width=\textwidth]{figures/real_robot_rollouts/output/frozen_vjepa21_scene.pdf}
    \caption{Matched visual reconstructions and recorded-action latent rollouts for the scene camera. Within each model, rows show the ground-truth (GT) video, independent frame reconstructions, and decoded autoregressive predictions.}
    \label{fig:robot-rollout-scene}
    \label{fig:robot-decode-panel}
\end{figure}
